\section{State of the art in plasma shape control at TCV}
\label{sec:sota_tcvshapectrl}
Despite TCV's redundant coil configuration and the broad range of plasma shapes it can support, the absence of a unified shape control system that fully integrates with existing shot design and preparation tools has remained a persistent limitation. Previous attempts at developing a shape controller for TCV have shed light on both the complexity of the task and the potential rewards brought by a successful solution. Early work reported in \cite{ariola_modern_1999,ariola_design_2002} introduced an $H_\infty$ controller aimed at managing key global shape parameters: triangularity, elongation, and vertical displacement. Although this approach represented a significant step forward, it relied heavily on a control system that is no longer in use at TCV, preventing its applicability in contemporary experiments. Another notable effort described in \cite{anand_novel_2017} proposed a novel shape control design but entirely replaced the pre-existing magnetic control architecture. This approach hindered seamless integration with the established framework and shot planning methodologies (see \Cref{sec:scd}), ultimately reducing its practical value.

More recently, a pioneering AI-based plasma shape controller developed in collaboration with Google DeepMind \cite{degrave_magnetic_2022} has shown remarkable capabilities during discharges. It has been developed to act on each phase of a shot, from the breakdown to the shape formation; it is currently being expanded to achieve more sophisticated plasma configurations. This innovative system highlighted the promise of advanced computational methods for achieving more agile and responsive control. However, its dependence on extensive computational resources (particularly for neural network tuning) has made it challenging to adapt and reconfigure between experiments.

The current approach is described in \cite{mele_design_2024}. It implements an isoflux control strategy to regulate the plasma boundary configuration. This approach maintains the plasma shape by enforcing equivalent poloidal flux values between a set of control points positioned along the last closed flux surface and a reference flux value measured at a specified \emph{boundary point}. The designation of this boundary point differs based on the plasma configuration: for \emph{limiter} plasmas, it corresponds to the limiter contact point, while for \emph{diverted} configurations, it is defined by the primary \emph{X-point} location. The control methodology varies according to the configuration type. In limiter scenarios, the controller enforces a zero orthogonal field component at the contact point to ensure the boundary intersects with the reference position. Diverted configurations, conversely, require simultaneous regulation of both radial and vertical field components to zero at the X-point location, thereby maintaining the desired magnetic topology. Furthermore, the study of some specific divertor configurations requires the control of other quantities such as poloidal flux at non-boundary points, flux difference between active and inactive X-points, and the position of inactive X-points; these too can be controlled by the shape controller.

Once a desired shape has been selected, the first step of the shape controller design procedure consists in determining a linearized model of the TCV plant around the prescribed equilibrium. This is quite a crucial and challenging task, usually solved by the suite of equilibrium reconstruction codes developed at SPC. The result of this design stage also yields a numerical representation of the plasma shape in terms of its LCFS, encoded as a set of control points.

The controller architecture implements a decoupling strategy for the controlled channels at low frequencies through singular value decomposition (SVD) of the steady-state plant dynamics. This decomposition is performed on the weighted static gain matrix of the system. To ensure both numerical feasibility and prevention of excessive control action, the controller design considers only singular values exceeding a specified threshold. This reduction in the problem dimensionality serves the dual purpose of maintaining computational feasibility while limiting control effort.

The architecture incorporates a first-order low-pass filter for each channel to attenuate oscillatory behavior. This design choice addresses the resonance phenomena observed in the Bode plots of the dominant singular modes across various TCV operational scenarios. Control of the selected dominant modes is achieved through a diagonal multiple-input multiple-output (MIMO) PID controller. When a state-space representation of the hybrid controller is available, two approaches exist for gain tuning: an analytical method based on assumed ideal channel decoupling with specified cutoff frequency and phase margin requirements per channel, or optimization-based techniques for refined parameter selection. The PID tuning methodology enables complete automation of the controller design process, upon the availability of a linearized plasma response model. The end user only has to specify the weighting matrices along with fundamental control parameters, such as the desired crossover frequency and phase margin.

This degree of automation proves particularly valuable during experiments, allowing the operators to be effective without extensive control theory expertise and ensuring an easy integration within the existing SCD framework.
